\emph{Översikt}
Den här modulen tillämpar klasser på två större exempel: en inköpslista
och ett enkelt banksystem.  Inköpslistan är en klass som har en
\emph{behållare} av varor som attribut och metoder som lägger till,
visar, bockar av och tar bort varor i den; vi frågar oss sedan om varje
vara ska vara en uppslagslista eller en egen klass, och väger
enkelhet mot struktur.  Banksystemet består av flera klasser som
samverkar --- person, adress, konto och medborgare --- och visar
\emph{komposition} (en klass har ett objekt av en annan som attribut)
och \emph{arv} (en klass specialiserar en annan) sida vid sida, så att
valet mellan dem blir tydligt.  Modulen avslutas med de designprinciper
exemplen praktiserar.

\emph{Lärandemål}
Efter denna modul ska studenten kunna:
% Lo09 (sammansatta datatyper), Lo02 (dela upp problem)
\begin{restatable}{lo}{AppLOContainer}\label{AppLOContainer}%
  Konstruera program där en klass har en behållare av andra objekt som
  attribut, och där metoderna söker i och uppdaterar behållaren.
\end{restatable}
% Lo09 (sammansatta datatyper), Lo06 (flexibla program)
\begin{restatable}{lo}{AppLORepresentation}\label{AppLORepresentation}%
  Välja mellan att representera data med en uppslagslista eller med en
  egen klass, och motivera valet utifrån programmets storlek och behov.
\end{restatable}
% Lo09 (sammansatta datatyper), Lo03 (dela upp program)
\begin{restatable}{lo}{AppLOComposition}\label{AppLOComposition}%
  Använda komposition för att bygga klasser av andra klasser.
\end{restatable}
% Lo01 (program utan kodupprepning), Lo09 (sammansatta datatyper)
\begin{restatable}{lo}{AppLOInheritance}\label{AppLOInheritance}%
  Använda arv för att specialisera en klass, och anropa
  föräldraklassens metoder med \mintinline{python}{super()}.
\end{restatable}
% Lo09 (sammansatta datatyper), Lo11 (granska program)
\begin{restatable}{lo}{AppLOChoice}\label{AppLOChoice}%
  Avgöra när komposition respektive arv passar bäst, och känna igen
  designprinciperna bakom valet i egen och andras kod.
\end{restatable}

\emph{Förkunskaper}
Modulen förutsätter klasser med attribut, metoder,
\mintinline{python}{__init__} och \mintinline{python}{__str__}, arv med
\mintinline{python}{super()} och egenskaper från föreläsningarna
\emph{Klasser och objekt} och \emph{Operatoröverlagring}, samt listor,
listbyggare (\foreignlanguage{english}{list comprehensions}) och
uppslagslistor från föreläsningarna om behållare.

\emph{Läsning}
Pythons egen handledning beskriver klasser och arv i kapitlet
\emph{Classes}.\footnote{\label{fn:pydoc-classes}%
  \url{https://docs.python.org/3/tutorial/classes.html} (hämtad
  2026-09-03).}
